\section{Introduction and main result}

A fiber is usually something one computes after the map has been chosen.
Here we start at the other end.  We choose the fiber first and ask whether a
Keller map can be built around it.

There is an immediate restriction.  A polynomial self-map of affine space
with nonzero constant Jacobian is \'etale, so every one of its fibers is a
finite \'etale scheme, possibly empty.  The question is whether this obvious
restriction is the only one: which finite \'etale algebras can occur as
\emph{entire} fibers of Keller maps?

The relevant condition is fullness.  Let \(K\) be a field and let
\[
 F:\A^m_K\longrightarrow\A^m_K
\]
be a Keller map.  Its geometric degree is
\[
 \gdeg(F)=[K(x_1,\ldots,x_m):K(F_1,\ldots,F_m)].
\]
For \(y\in\A^m(K)\), we call \(F^{-1}(y)\) a \emph{full Keller fiber} if
\[
 \operatorname{rank}_K F^{-1}(y)=\gdeg(F).
\]
Thus the fiber contains every generic inverse sheet; no boundary calculation
is needed once the two ranks have been identified.  Fullness is what turns a
special fiber from a partial coincidence into a faithful realization of the
generic covering degree.

Our main theorem is the characteristic-zero statement proved end to end in
Lean.

\begin{theorem}[Characteristic-zero polynomial-presentation theorem]
\label{thm:main-intro}
Let \(K\) have characteristic zero, and let \(P\in K[T]\) be separable of
degree \(N\geq3\).  There are a polynomial map
\[
 F_P:\A^3_K\longrightarrow\A^3_K
\]
and a target \(y_P\in\A^3(K)\) such that
\[
 \det DF_P=1,\qquad
 \gdeg(F_P)=N,\qquad
 F_P^{-1}(y_P)\simeq\Spec K[T]/(P).
\]
The fiber isomorphism is natural on commutative \(K\)-algebras, every
coordinate of \(F_P\) has degree at most \(6N+2\), and the displayed fiber is
finite \'etale of rank \(N\).
\end{theorem}

The theorem says more than that the same number of points can be arranged.
It realizes the quotient algebra itself, naturally over every commutative
test algebra, and the fiber has exactly the geometric degree of the map.
Equivalently, separability means \(\gcd(P,P')=1\); this is the
derivative-squarefree condition intended throughout the paper.
Theorem~\ref{thm:polynomial-realization} gives the proof.  Its exact
paper-facing Lean declaration is
\texttt{paperPolynomialPresentation\_pageOne}.

The construction appeared only after reversing the usual direction of
thought.  Instead of guessing three source coordinates and then trying to
understand their fibers, we begin with the one-variable equation we want to
solve in the inverse direction.  Its derivative will do two jobs: it is the
unit that reconstructs a source point from a root, and it is the factor that
disappears from the Jacobian.  To insert \(P\), we translate it and use
\[
 G(S)=P(a+S)-P(a)
\]
as the seed of this inverse equation.

The explicit formulas require the translation \(a\) to satisfy two
nonvanishing conditions.  If \(a\) is kept as supplied data, the same idea
works under broader characteristic hypotheses.

\begin{proposition}[Supplied-presentation version]\label{prop:supplied-intro}
Assume that \(\operatorname{char}K\neq2\).  Let \(P\in K[T]\) be separable
of degree \(N\geq3\), and let \(a\in K\) satisfy
\[
 P^{[1]}(a)P^{[3]}(a)\neq0,
\]
where \(P^{[j]}\) denotes the \(j\)-th Hasse derivative.  There are an
explicit polynomial map
\[
 \widetilde F_{P,a}:\A^3_K\longrightarrow\A^3_K
\]
and an explicit target \(y_{P,a}\in\A^3(K)\) such that
\[
 \det D\widetilde F_{P,a}=1,\qquad
 \gdeg(\widetilde F_{P,a})=N,\qquad
 \widetilde F_{P,a}^{-1}(y_{P,a})
 \simeq\Spec K[T]/(P).
\]
The fiber isomorphism is natural on commutative \(K\)-algebras, every
coordinate has degree at most \(6N+2\), and the supplied
\((P,a)\)-construction commutes with extension of the ground field.
\end{proposition}

Proposition~\ref{prop:supplied-polynomial-realization} gives the formulas and
proof, while Proposition~\ref{prop:base-change} gives scalar-extension
compatibility.  This characteristic-\(\ne2\) proposition is not yet
formalized end to end.

The pleasing feature is that the entire three-variable map is governed by
the compact scalar inverse equation
\[
 E_{\Pi,B,C}(S)=0
\]
whose derivative reconstructs the source and cancels in the Jacobian.  Over
the generic target, the marked root \(S\) generates the actual source field.
Because \(E\) is linear in one target coordinate, its irreducibility and
degree are then transparent.  At the distinguished target the inverse
equation becomes \(E(S)=P(a+S)\); squarefreeness removes the derivative
localization and leaves the literal fiber scheme.  In one line, fullness is
\[
 \operatorname{rank}_K\widetilde F_{P,a}^{-1}(y_{P,a})
 =N=\gdeg(\widetilde F_{P,a}).
\]

This is already the polynomial-presentation theorem.  To pass from a
polynomial to an abstract algebra, one more classical observation completes
the picture: every finite \'etale algebra over an infinite field is
monogenic.  In characteristic zero an admissible translation always exists.
Monogenicity therefore gives the paper's principal abstract corollary:

\begin{corollary}\label{cor:main-intro}
If \(K\) has characteristic zero, every finite \'etale \(K\)-algebra of rank
at least three occurs as a full Keller fiber of a Jacobian-one polynomial
map \(\A^3_K\to\A^3_K\).
\end{corollary}

The word ``explicit'' applies to
Proposition~\ref{prop:supplied-intro}, where both the polynomial presentation
and admissible translation are supplied.  In the headline theorem the
admissible translation is chosen automatically but noncanonically.  Starting
from an abstract finite \'etale algebra, the primitive element is also chosen
noncanonically; the corresponding Lean construction is noncomputable and no
functorial choice is claimed.

There is one small rank that does not join the pattern.  In characteristic
zero, the construction realizes every rank \(N\geq3\).  The classical
Galois case of the Jacobian theorem, due to Campbell, Razar, and Wright,
excludes degree two.  Thus the possible nonzero ranks are
\[
 1,3,4,\ldots.
\]
This classification is a corollary with an external input, not part of the
formalized construction.

The paper follows the route by which the construction becomes visible.
Section~\ref{sec:gauge} begins with the inverse equation and pulls it back to
a polynomial map.  Section~\ref{sec:reconstruction} shows that a marked root
and a source point are the same data, then identifies the actual
function-field degree.  Section~\ref{sec:prescribing} inserts a prescribed
separable polynomial and then an abstract finite \'etale algebra.
Section~\ref{sec:example} lets the construction run on a concrete arithmetic
fiber.  The appendices keep the coefficient expansion and Lean
correspondence visible without interrupting the main story.

The formal boundary is exact.  The characteristic-zero headline theorem and
abstract corollary are formalized end to end.  The separate
characteristic-\(\ne2\) supplied-presentation proposition is proved directly
here, but its current Lean files still assume \texttt{CharZero}.  Exact
nonproperness, symmetric monodromy, Hilbertian specialization, and stable
compositional atomicity belong to separate geometric and monodromy results;
none is used in this paper.
